\section{Related Work}

\subsection{Proficiency resources and grammatical representation}

The CEFR provides proficiency scales across language activities rather than a
machine-readable inventory of grammatical skills \citep{coe-2020-cefr}.  The
English Profile programme complements these scales with empirical descriptions
of grammatical and lexical competence.  The EGP draws on learner-corpus
frequency, accuracy, and dispersion to associate grammar competence statements
with CEFR levels \citep{okeeffe-mark-2017}.  Its level assignments have also
been revisited quantitatively using NLP methods
\citep{verratti-souto-etal-2025}.  The present work does not reassess CEFR
levels.  It instead asks how the grammatical content of a bounded EGP subset
can be represented without converting prose descriptors directly into KCs or
silently completing information that the source does not provide.

This distinction separates a \emph{curriculum source} from a \emph{domain
model}.  Multiple descriptors can support one grammatical configuration, one
descriptor can explicitly support alternatives, and many descriptors do not
fully determine all properties required by a closed ontology.  \system retains
those cases as partial or out of scope rather than forcing them into an exact
cell.

\subsection{Knowledge components, Q-matrices, and model selection}

Classical Bayesian Knowledge Tracing (BKT) models mastery as a latent binary
skill state updated after each opportunity \citep{corbett-anderson-1995}; Deep
Knowledge Tracing replaces explicit Bayesian state transitions with a neural
sequence model \citep{piech-etal-2015}.  Broader surveys show that modern KT
architectures differ substantially in their state representations and use of
item information \citep{abdelrahman-etal-2023}.  All nevertheless require a
definition of the content whose learning is being traced.

KCs are not necessarily natural atoms waiting to be annotated.  The
Knowledge--Learning--Instruction framework treats them as explanatory units
linking tasks, learning processes, and instruction \citep{koedinger-etal-2012}.
Learning Factors Analysis evaluates and refines a tutor's cognitive model using
student data and explicit complexity trade-offs \citep{cen-etal-2006}.  Related
factor models likewise express performance through skill difficulty and
practice terms \citep{pavlik-etal-2009}.  These approaches motivate treating a
KC inventory as a hypothesis rather than administrative metadata.

In cognitive diagnosis, the Q-matrix specifies which latent attributes are
required by each item.  Expert construction is subjective, and
misspecification can affect learner classification; empirical validation and
refinement methods therefore examine the adequacy of proposed entries
\citep{de-la-torre-chiu-2016,pardos-dadu-2018}.  \system intervenes before human
response data are available.  It tests the structural support, scope, and
identifiability of declared candidate rules, freezes a policy, and leaves later
learner-data validation as a separate stage.

\subsection{Knowledge tracing in tutoring dialogue}

Open-ended tutoring dialogue does not arrive with question identifiers,
correctness labels, or a fixed KC set.  \citet{scarlatos-etal-2025-dialogue}
frame tutor--student dialogue as a sequence of formative assessments and use
LLMs to annotate turn-level KCs and correctness for CoMTA and MathDial.  Their
pipeline retrieves mathematics standards as candidate KCs and evaluates the
automatic annotations against expert labels.  \citet{huang-etal-2026-interpretable}
extend conversational KT by explicitly modelling student ability and task
difficulty through an LLM and item-response-theoretic representation.

\system addresses a complementary, upstream problem.  Rather than assigning a
turn to an already available skill standard, it constructs and compares the
skill representation itself.  TSCCv2 provides a plausible eventual test bed
because it includes learner language and grammatical corrections
\citep{caines-etal-2022-teacher}; however, applying the present ontology to
TSCCv2 would still require a separate definition of assessment opportunities,
context-sensitive correctness, and multi-phenomenon turn annotation.

\subsection{Exercise generation for language learning}

Prior work has generated educational questions conditioned on learner history
or target difficulty \citep{srivastava-goodman-2021-question,cui-sachan-2023-adaptive}.
Grammar-oriented systems generate reading, multiple-choice, or gap-filling
exercises from passages and examples
\citep{chan-etal-2022-agree,bitew-etal-2023-learning}.  Such work demonstrates
both the cost of manual authoring and the importance of correctness,
naturalness, and answer uniqueness.

The present item generator deliberately makes a different trade-off.  It uses
controlled transformations with fixed lexical frames so that the relationship
among a canonical cell, a concrete surface derivation, and a later KC
projection remains inspectable.  This reduces lexical and discourse diversity
but avoids using unconstrained generation to hide changes in the experimental
content model.

\subsection{Synthetic learners and controlled evaluation}

Synthetic data can test whether a KT implementation recovers known patterns or
responds to controlled changes \citep{pagonis-etal-2024}.  It is also easy to
overinterpret.  Recent evaluation of LLM-simulated students finds substantial
behavioural and cognitive limitations even when linguistic output appears
plausible \citep{scarlatos-etal-2026-simulated}.  \system therefore uses a
transparent numerical structural oracle rather than an LLM role-play student,
and calls its outputs an experimental world rather than human learner data.
The oracle is useful for testing leakage, component reuse, cold-start behaviour,
and compositional transfer; it cannot establish the cognitive validity or
acquisition order of grammar KCs.
